\documentclass[12pt,a4paper]{amsbook}

\input{preamble}

\title[Longform template example]{Longform template example}
\author{Author}
\date{\today}

\begin{document}

\frontmatter

\maketitle

\tableofcontents

\chapter*{Preface}

This example demonstrates the longform template: front matter, chapters,
sections, theorem-like environments, cross-references, citations, and a
bibliography.

\mainmatter

\chapter{First chapter}

\section{Basic structure}

This is ordinary mathematical text in a chaptered document. We cite
\cite{har77} and refer to \Cref{theorem:sample}.

\begin{definition}\label{definition:sample}
A \emph{sample object} is an object introduced to demonstrate the formatting of
long-form mathematical writing.
\end{definition}

\begin{theorem}\label{theorem:sample}
Let \(X\) be a sample object. Then \(X\) satisfies the following properties:
\begin{enumerate}[label=(\roman*)]
  \item \label{item:longform-first} \(X\) has a first property.
  \item \label{item:longform-second} \(X\) has a second property.
\end{enumerate}
\end{theorem}

\begin{proof}
This is a sample proof. The displayed equation
\[
  \chi(\mathscr{F}) =
  \sum_i (-1)^i h^i(X,\mathscr{F})
\]
checks the display style in a long-form document.
\end{proof}

The first property is \Cref{item:longform-first}.

\begin{proposition}\label{proposition:sample}
The sample object of \Cref{definition:sample} also satisfies a proposition.
\end{proposition}

\begin{remark}\label{remark:sample}
This is a remark in the longform template.
\end{remark}

\chapter{Second chapter}

\section{A diagram}

The template supports commutative diagrams:
\[
\begin{tikzcd}
A \arrow[r] \arrow[d] & B \arrow[d] \\
C \arrow[r] & D .
\end{tikzcd}
\]

\section{References inside the document}

We have used \Cref{definition:sample,theorem:sample,proposition:sample} and
the citation \cite{har77}.

\backmatter

\printbibliography

\end{document}
